\section{引言}\label{intro}

若干年前，A. Kapustin 与作者在文献 \cite{KW} 中表明，几何朗兰兹对应\index{geometric-langlands-correspondence@几何朗兰兹对应}（导论可参见例如文献 \cite{frenkel}）可以用四维 $\N=4$ 超杨-米尔斯理论\index{super-yang-mills-theory@超杨-米尔斯理论}的 $S$-对偶性来表述. 基本思想是在紧致化\index{compactification@紧致化}到二维的背景下，考察该理论的边界条件\index{boundary-condition@边界条件}、$S$-对偶性对这些边界条件的作用，以及超对称\index{supersymmetry@超对称}线算符\index{line-operator@线算符}在边界附近的行为. 此后，人们作出了若干推广. 通过将面算符\index{surface-operator@面算符}以及位于其上的线算符纳入分析，这一理论已被推广到包括温和分歧\index{tame-ramification@温和分歧} \cite{GW} 和野分歧\index{wild-ramification@野分歧} \cite{Wild} 的情形. 文献 \cite{FW} 探讨了支撑在奇点处的膜所涉及的几何朗兰兹对偶性中最简单的问题. 文献 \cite{GaiW,GaiW2} 对半 BPS 边界条件、畴壁及其在 $S$-对偶性下的行为所作的研究，部分动机来自它们在几何朗兰兹中的潜在应用. 几何朗兰兹的规范理论\index{gauge-theory@规范理论}方法中的某些要素，也与通过规范理论理解 Khovanov 同调有关 \cite{WittenKhov}. 关于规范理论与几何朗兰兹某些方面的非正式说明，参见文献 \cite{EW1,EW2,EW3}.

本文旨在介绍与这一主题有关的一些补充内容. 我们只考虑无分歧\index{unramified@无分歧}情形，因此不引入面算符. 我们的主要目标是具体说明（不假设电磁对偶性\index{electromagnetic-duality@电磁对偶性}），为什么希格斯丛\index{higgs-bundle@希格斯丛}模空间\index{moduli-space@模空间} $\MH$ 上某些简单的 $A$-膜，正如对偶性所预言的那样，是“磁本征膜\index{magnetic-eigenbrane@磁本征膜}”. 如文献 \cite{KW} 所述，这里所考虑的是支撑在希钦纤维化\index{hitchin-fibration@希钦纤维化}的一条纤维上、与零维膜对偶的秩为 1 的膜. 为了直接证明这些膜是磁本征膜，我们首先需要回顾希格斯丛和希钦纤维化的一些事实.

虽然我们重述了一些细节，以使本文更加自足，但读者可能仍需熟悉前文 \cite{KW} 的部分内容. 我们的记号大体沿用该文，只作少量修改. 特别地，与文献 \cite{KW} 一样，我们用 $^L\neg G$ 表示紧李群 $G$ 的朗兰兹对偶群或 GNO 对偶群（在上面引用的某些论文中，对偶群记为 $G^\vee$）.

作为准备，我们在第 \ref{compacthitchin} 节和第 \ref{dualtori} 节回顾希格斯丛与希钦纤维化的一些基本事实. 这些内容都不是新的，但其中有些可能不太容易查到.（在这两节中，我们会解释一些有用的细节，尽管就我们的应用而言并非严格必需.）在第 \ref{thooftheckeop} 节，我们遵循文献 \cite{KW}，介绍“本征膜\index{eigenbrane@本征膜}”的概念，描述 $\MH(^L\neg G)$ 上的电本征膜\index{electric-eigenbrane@电本征膜}，并说明 $S$-对偶性对 $\MH(G)$ 上的磁本征膜所作的预言. 本文的其余部分致力于直接证明，所讨论的膜确实是磁本征膜. 我们只在最简单的情形中这样做，也就是对那些使这一论证最为直接的规范群和 't Hooft 算符进行讨论.

\section{紧致化与希钦模空间}
\label{compacthitchin}

\subsection{预备知识}\label{prelims}

从四维规范理论\index{gauge-theory@规范理论}导出几何朗兰兹对应\index{geometric-langlands-correspondence@几何朗兰兹对应}之所以自然，根本原因是 $2+2=4$. 在通常的数学表述中，几何朗兰兹对应将两类范畴联系起来，它们分别与黎曼曲面 $C$ 的基本群的表示以及其上的全纯向量丛\index{vector-bundle@向量丛}相关. 从物理观点看，基本群的一个表示或一个全纯向量丛由 $C$ 上的一个{\it 规范场\index{gauge-field@规范场}}描述. 因此，我们应当预期要研究某种规范理论. 但是，我们必须考虑到这样一个事实：几何朗兰兹对应是关于与 $C$ 相关的{\it 范畴}的陈述，而不是关于数或向量空间的陈述.

一个 $d$ 维量子场论给一个 $d$ 维流形 $M_d$ 赋予一个数，即配分函数或路径积分的值. 对于一个 $d-1$ 维流形 $M_{d-1}$，它赋予一个向量空间，即在 $M_{d-1}$ 上对该理论进行量子化所得的量子态空间 $\H_d$. 这一层级中的下一步稍微不那么为人熟知（参见例如文献 \cite{Freed,BD,
Kapustin}）. 一个 $d$ 维量子场论给一个维数为 $d-2$ 的流形 $M_{d-2}$ 赋予一个{\it 范畴}. 例如，一个 $d=2$ 的场论给一个零维流形，也就是一个点，赋予一个范畴. 由于任意两点都是同构的，这仅仅意味着一个二维场论确定了一个范畴，即边界条件\index{boundary-condition@边界条件}的范畴.（最熟悉的例子可能是二维拓扑扭曲 $\sigma$-模型中的 $A$-膜和 $B$-膜的范畴.）

因此，如果我们希望一个量子场论给任意二维流形 $C$ 赋予一个范畴，就必须从一个 $2+2=4$ 维的量子场论出发. 此外，如果这个范畴应当与 $C$ 上的规范场相关，我们就应从四维规范场出发. 最后，如果我们希望找到一种对偶性，将一个以某种方式与规范群 $G$ 相关的范畴，同一个以某种方式与朗兰兹或 GNO 对偶群 $^L\neg G$ 相关的范畴联系起来，那么就应当从一种具有交换这两个群的对偶性的四维理论出发. 具备所需性质的规范理论就是四维 $\N=4$ 超杨-米尔斯理论\index{super-yang-mills-theory@超杨-米尔斯理论}.\transnote{这里的 $2+2$ 指四维时空分解为紧致化所用的黎曼曲面 $C$（2 维）与剩余二维流形 $\Sigma$；$d$ 维理论对 $d-2$ 维流形赋予范畴，故要得到与 $C$ 有关的范畴必须取 $d=4$，同时剩余方向恰为二维，才能约化出 $\Sigma$ 上的二维 $\sigma$-模型. 原文 §2.1.（来源：A 级，arXiv:1506.04293 §2.1；A 级，arXiv:hep-th/0604151（Kapustin–Witten 2006）摘要，2026-11-19 核对.）}

为了应用于几何朗兰兹，我们对 $\N=4$ 理论作拓扑扭曲，从而在形式上得到一个四维拓扑场论.\footnote{\label{partial}这很可能只是一个部分拓扑场论. 一般而言，它的配分函数是否良定义尚不清楚；要定义配分函数，就必须处理非紧零模空间\index{moduli-space@模空间}上的积分.（在一个四维流形 $M$ 上，作用量的经典极小值所构成的空间，是模去共轭的同态 $\pi_1(M)\to G_\C$ 的空间，一般是非紧的. 这个空间的非紧性，会给拓扑不变性的证明，以及拓扑场论应有的切割与粘合性质的证明带来困难.）然而，对于那些不涉及场空间非紧性的可观测量，其定义是良好的. 特别地，存在良定义的膜范畴，而与几何朗兰兹有关的其他可观测量也是良定义的. 从技术上说，由于 $\MH$ 非紧，可以对膜在无穷远处的行为施加不同条件；为了在细节上与几何朗兰兹的标准数学表述一致，需要作出正确的选择.} 关于这种扭曲，以及它如何产生由 $S$-对偶性交换的一对拓扑场论，即 $A$-模型和 $B$-模型的细节，参见文献 \cite{KW}.\transnote{所谓“部分”拓扑场论，指配分函数一般并不良定义：四维流形 $M$ 上经典真空的模空间是模去共轭的 $\pi_1(M)\to G_{\mathbb C}$ 的同态空间，一般非紧，故路径积分与切割、粘合性质都难以证明；但膜范畴等不涉及该非紧性的可观测量仍良定义，这正是本篇只需膜范畴即可推进的原因. 原文 §2.1 作者原注.（来源：A 级，arXiv:1506.04293 §2.1 原注；A 级，arXiv:hep-th/0604151，2026-11-19 核对.）}

只需考虑在黎曼曲面 $C$ 上从四维紧致化\index{compactification@紧致化}到二维，就引入了 $C$. 因此，我们考察四维流形 $\R^2\times C$ 上的四维 $\N=4$ 理论，或者更一般地，考察 $\Sigma\times C$ 上的理论，其中 $\Sigma$ 是一个假定远大于 $C$ 的二维流形. $\Sigma\times C$ 上的拓扑场论（如脚注所述，是部分拓扑场论）在 $C$ 上紧致化之后，约化为 $\Sigma$ 上的一个拓扑场论.

在研究 $\Sigma$ 上的这一理论时，我们总可以假定 $\Sigma$ 远大于 $C$，并在 $\Sigma$ 上的大距离尺度下工作. 这意味着，对于许多目的，我们可以用 $\Sigma$ 上的二维理论来代替 $\Sigma\times C$ 上的完整四维理论 \cite{BJSV,HMS}. 这个理论是一个超对称\index{supersymmetry@超对称} $\sigma$-模型，其靶空间是在 $C$ 上紧致化时产生的经典超对称真空的空间. 在远离部分规范对称性恢复的奇点\footnote{在这一陈述中，我们假设 $C$ 的亏格 $g_C$ 至少为 2，或者在亏格为 0 或 1 的情形中引入分歧. 当亏格为 0 或 1 且没有分歧时，希钦方程\index{hitchin-equation@希钦方程}的解具有阿贝尔结构群，需要更仔细地表述. 我们还假设 $G$ 是单群或半单群. 对于 $G=U(1)$ 或 $U(N)$，情形有所不同，因为 $G$ 总有一个未破缺的 $U(1)$ 子群. 关于这一点的讨论，参见第 \ref{elmag} 节.}处，这个空间是 $C$ 上结构群为 $G$ 的“希格斯丛\index{higgs-bundle@希格斯丛}”的希钦超凯勒\index{hyperkahler@超凯勒}模空间. 我们将这个空间记为 $\MH$；如果希望明确指定规范群 $G$ 和进行紧致化的黎曼曲面 $C$，则记为 $\MH(G,C)$.\transnote{“希格斯丛”借自粒子物理而含义不同：这里的希格斯场 $\varphi$ 是 $C$ 上取值于 $\ade$ 的一次形式，并非标量场；它来自四维 $\mathcal N=4$ 理论经拓扑扭曲后沿 $C$ 约化所保留的数据. 该模空间带 $I,J,K$ 三个复结构，故称超凯勒流形. 原文 §2.1.（来源：B 级 nLab “Higgs bundle” \url{https://ncatlab.org/nlab/show/Higgs+bundle}，2026-11-19 访问；与 A 级 arXiv:1506.04293 §2.1–§2.2 交叉核对.）}

虽然用靶空间为 $\MH$ 的 $\sigma$-模型来描述十分有用，但这种描述实际上只在一般位置成立. 这是因为，$\MH$ 的某些奇点处会恢复部分规范对称性，必须引入额外的自由度，才能得到良好的描述. 究竟需要引入多少额外自由度，取决于所提出的问题. 不作任何近似、对所有目的都有效的普适描述，是其背后的 $\Sigma\times C$ 上的四维规范理论，无需尝试将其约化为有效二维理论. 如果想把这个完整的四维理论描述成一个二维 $\sigma$-模型，就应当称它为如下二维超对称 $\sigma$-模型：靶空间为 $C$ 上全部规范场的空间 $\Y$（更准确地说，如我们稍后将讨论的，为了包括希格斯场\index{higgs-field@希格斯场}，应取对应的余切丛\index{cotangent-bundle@余切丛} $\CW=T^*\Y$），规范群则为从 $C$ 到有限维群 $G$ 的所有映射组成的无限维群 $\G$. 如果对这一 $\sigma$-模型作拓扑扭曲，得到 $A$-模型或 $B$-模型，那么 $\G$ 实际上就被它的复化所代替，即从 $C$ 到 $G$ 的复化 $G_\C$ 的映射所构成的群 $\G_\C$.

为了将这一图景与几何朗兰兹中的标准陈述联系起来，我们应当具体考察对偶性中的 $A$-模型一侧或 $B$-模型一侧. 为简洁起见，这里只考虑 $A$-模型一侧；类似的陈述也适用于 $B$-模型. 代数几何学家用 $C$ 上所有全纯 $G$-丛的“叠\index{stack@叠}”来表述 $A$-模型，这些丛不必稳定.\footnote{因此，根据文献 \cite{BeD}，在对偶性的一侧，必须考虑 $G$-丛的叠上的 $\mathcal D$-模（微分算符层上的模），而不仅是稳定 $G$-丛模空间上的模. $A$-模型与 $\mathcal D$-模的关系，依赖于希格斯丛模空间可以由一个余切丛近似这一事实，文献 \cite{KW} 对此作了解释.} 然而，根据 Atiyah 和 Bott \cite{AB}，这个叠的一个具体模型，就是由无限维群 $\G_\C$ 作用的无限维空间 $\Y$. 因此，物理学家对“完整描述涉及整个‘叠’”这一说法的理解，仅仅是说，完整描述涉及一个以整个无限维空间 $\Y$ 为靶空间（并带有 $\G_\C$ 的作用）的超对称 $\sigma$-模型，换言之，涉及完整的四维规范理论.

\subsection{希钦方程}\label{hitcheq}

不过，在本文中，我们主要考虑只需用靶空间为 $\MH$ 的 $\sigma$-模型就足以处理的问题. $\MH$ 由一组熟悉的方程定义，称为希钦方程.\transnote{这组方程由 Hitchin 在 1987 年研究黎曼曲面上的自对偶方程时引入；把 $A$ 与 $\phi$ 合并为复值联络后，条件 (2.2) 就是四维自对偶方程沿 $C$ 约化的形式. 注意 $\phi$ 取值于 $\ade$ 且是一次形式，这是它与粒子物理中标量希格斯场的本质差别. 原文 §2.2 式 (2.1)–(2.2).（来源：A 级，Hitchin, Proc. London Math. Soc. (3) 55 (1987) 59–126（本文文献 [13]）；B 级 nLab “Higgs bundle”，2026-11-19 访问.）} 这些是关于一对场 $A,\phi$ 的方程，其中 $A$ 是二维曲面 $C$ 上的一个 $G$-丛 $E$ 上的联络，而 $\phi$ 是 $C$ 上取值于伴随表示的一个 1-形式，也就是说，它取值于与 $E$ 相关联的伴随丛 $\ade$. 在写出这些方程时，假定 $C$ 是一个已经选定复结构\index{complex-structure@复结构}的黎曼曲面. 在 $C$ 上引入局部复坐标 $z$，这些方程可以写成
\begin{align}
\label{secondway} F_{z\bar z}-[\phi_z,\phi_{\bar z}]&=0\nonumber
\\    D_{\bar z}\phi_z=D_z\phi_{\bar z}&=0. \end{align}
这里 $D_{\bar z}=\partial_{\bar z}+[A_{\bar z},\cdot],$ $D_z=\partial_z+[A_z,\cdot]$.

另一种做法是将 $A$ 与 $\phi$ 合并成复值联络 $\CA=A+i\phi$. 我们将其视为 $G_\C$-丛 $E_\C\to C$ 上的联络，这个丛由 $E\to C$ 复化而来. 它的结构群为 $G_\C$，曲率为 $\F=d \A+\A\wedge \A$. $\F$ 的实部和虚部分别为 ${\rm Re}\,\CF=F-\phi\wedge\phi$、${\rm Im}\,\CF=D\phi$. 希钦方程等价于
\begin{align}\label{firstview}
 \F & = 0 \nonumber \\
                     D\star \phi&=0. \end{align}
这里 $D\star\phi=D_{\bar z}\phi_z+D_z\phi_{\bar z}$.

希钦方程的这两种写法，在这些方程的解的模空间 $\MH$ 上给出两个不同的复结构. 从式 (\ref{firstview}) 出发，我们看到，方程 $\F=0$ 在 $\A$ 的复值规范变换下不变，而模去这样的规范变换后，解空间就是模去共轭的同态 $\pi_1(C)\to G_\C$ 的等价类空间. 特别地，方程 $\F=0$ 是拓扑不变的；它实际上不依赖于 $C$ 上的复结构. 一个经典定理 \cite{Corlette} 表明，在一个温和且一般成立的半稳定\index{semistability@半稳定}性条件下，\footnote{这一条件是说，如果平坦联络\index{flat-connection@平坦联络} $\A$ 的单值变换\index{monodromy@单值变换}能够同时共轭成分块三角形式 $\begin{pmatrix} * & * \cr 0 & * \end{pmatrix}$，那么事实上它们可以共轭成分块对角形式 $\begin{pmatrix} * & 0 \cr 0 & * \end{pmatrix}$. 这个条件一般会满足，因为（对于亏格 $\geq 2$ 的 $C$）一般而言，$C$ 上一个复平坦联络的单值变换不能共轭成分块三角形式.} 按取值于 $G_\C$ 的规范变换对 $C$ 上的复平坦联络进行分类，等价于施加“矩映射\index{moment-map@矩映射}”方程 $D\star\phi=0$，然后只模去取值于 $G$ 的规范变换.

因此，从一个角度看，$\MH$ 可以解释为半稳定同态 $\pi_1(C)\to G_\C$ 的模空间. 由于 $G_\C$ 本身是一个复流形，这一解释赋予 $\MH$ 一个复结构. Hitchin 将这个复结构称为 $J$ \cite{Hitchin}.\transnote{复结构 $J$ 来自“半稳定同态 $\pi_1(C)\to G_{\mathbb C}$ 的模空间”这一描述，复结构 $I$ 则来自希格斯丛；两种描述之所以一致，依据的是非阿贝尔 Hodge 对应，而本文所用的半稳定性条件正是它的前提. 原文 §2.2.（来源：A 级，arXiv:1506.04293 §2.2；B 级 nLab “Higgs bundle” 中 nonabelian Hodge theory 一节，2026-11-19 访问.）}

从另一个角度看，我们将 $A_{\bar z}$ 和 $\phi_z$ 视为全纯变量（并将它们的复共轭 $A_z$ 和 $\phi_{\bar z}$ 视为反全纯变量）. 希钦方程 (\ref{secondway}) 因而分解为一个全纯方程 $D_{\bar z}\phi_z=0$（它蕴含复共轭方程 $D_z\phi_{\bar z}=0$）和“矩映射”方程 $F-\phi\wedge\phi=0$. 为了解释这个全纯方程，我们首先注意到，对于丛 $E\to C$ 上任意取值于 $G$ 的联络 $A$，一旦选定 $C$ 上的复结构，相应的 $\bar\partial$ 算符 $D_{\bar z}$ 就赋予 $E$ 一个全纯结构. 方程 $D_{\bar z}\phi_z=0$ 是说，在这个复结构中，$\phi_z$ 可以视为 $K\otimes \ade$ 的一个全纯截面，其中 $K$ 是 $E$ 的典范丛\index{canonical-bundle@典范丛}. 按照 Hitchin 的术语，一个希格斯丛是一对数据：一个全纯 $G_\C$-丛 $E\to C$，以及一个全纯“希格斯场” $\varphi=\phi_z d z\in H^0(C,K\otimes \ade)$；换言之，就是方程 $D_{\bar z}\phi_z=0$ 的一个解. 我们将这样的一对 $(E,\varphi)$ 称为希钦对\index{hitchin-pair@希钦对}.

取值于 $G_\C$ 的规范变换自然作用在满足 $D_{\bar z}\phi_z=0$ 的对 $A_{\bar z}$、$\phi_z$ 所构成的空间上. 不过，Hitchin 证明，在同样施加一个温和的半稳定性条件\footnote{\label{semicon}为简单起见，我们只在 $G=SU(N)$ 的情形中陈述这个条件. 此时，如果任意子丛的第一陈类\index{chern-class@第一陈类}都非正，则全纯 $G_\C$-丛 $E_\C\to C$ 是半稳定的. 如果任意 $\varphi$-不变子丛的第一陈类都非正，则希格斯丛 $(E,\varphi)$ 是半稳定的.（$\varphi$-不变性的概念如下. 由于 $\varphi$ 是 $K\otimes \ade$ 的一个截面，它可以自然地理解为一个映射 $E\to K\otimes E$. 选取 $K$ 的一个局部平凡化\index{trivialization@平凡化}，就得到一个映射 $E\to E$；如果 $E$ 的一个子丛在这个映射下不变，就称它是 $\varphi$-不变的. 这一判据不依赖于 $K$ 的局部平凡化的选择.）特别地，如果 $E$ 是半稳定的，那么对任意 $\varphi\in H^0(C,K\otimes\ade)$，$(E,\varphi)$ 都是半稳定的. 在这些陈述中，若将“半稳定”换成“稳定”，只需将“非正”换成“负”. 在术语方面，应当指出，通常所谓某类稳定对象的模空间，一般定义为参数化该类型的稳定或半稳定对象. 这就是为什么人们把 $\MH$ 称为稳定希格斯丛的模空间，尽管一般而言，它所参数化的某些对象仅仅是半稳定的.}之后，将这些对组成的空间模去取值于 $G_\C$ 的规范变换，等价于施加“矩映射”方程 $D\star\phi=0$，然后只模去取值于 $G$ 的规范变换.

因此，从这个观点看，$\MH$ 又获得了另一种解释，即稳定希格斯丛的模空间. 这种描述使另一个复结构显而易见，Hitchin 将其称为 $I$. 此外，Hitchin 还证明，复结构 $I$ 和 $J$ 是 $\MH$ 上一个超凯勒结构的组成部分；特别地，四元数单位 $I$、$J$ 和 $IJ=K$ 作用在 $\MH$ 的任意光滑点处的切空间上.\transnote{$I,J,K$ 满足四元数关系 $IJ=K$、$I^2=J^2=K^2=-1$，故在每点切空间上给出四元数模结构；超凯勒流形上还有以单位球面参数化的复结构族，本篇只用到其中的 $I$ 与 $J$. 原文 §2.2.（来源：A 级，arXiv:1506.04293 §2.2；A 级，Hitchin, Proc. London Math. Soc. 55 (1987)（本文文献 [13]），2026-11-19 核对.）}

这一图景中与几何朗兰兹有关的许多方面，已在文献 \cite{KW} 中作了说明，这里不再重复. 关于 $\MH$ 是超凯勒的这一陈述，我们只指出，它可以构造成一个超凯勒商. 由对 $A,\phi$ 组成的空间 $\CW$ 可以视为一个无限维超凯勒流形，具有平坦超凯勒度量
\begin{equation}
\label{igob}
ds^2=-{1\over
2\pi}\int_C|d^2z|\,\Tr\bigl(\delta A_z\otimes \delta A_{\bar
z}+\delta \phi_z\otimes
\delta \phi_{\bar z}\bigr),~~~|d^2z|=i\,dz\,d\bar z.
\end{equation}
取值于 $G$ 的规范变换群 $\G$ 作用在 $\CW$ 上，并保持其超凯勒结构. 希钦方程断言，$\G$ 在 $\CW$ 上的作用所对应的超凯勒矩映射为零. 因此，将希钦方程的解模去 $\G$ 所得的空间 $\MH$，可以解释为超凯勒商 $\CW///\G$. 由此，它自然带有一个超凯勒结构.\transnote{超凯勒商要求两个矩映射同时为零：实矩映射对应 $D\star\phi=0$，复矩映射对应 $F-\phi\wedge\phi=0$，合起来正是式 (2.1) 的两组方程；文中“超凯勒矩映射为零”即指此. 原文 §2.2 式 (2.1)–(2.4).（来源：A 级，arXiv:1506.04293 §2.2；A 级，arXiv:hep-th/0604151，2026-11-19 核对.）}

这里，我们将着重讨论复结构 $I$ 中这一图景的若干方面，它们在文献 \cite{KW} 中没有得到解释，却对更深入地理解几何朗兰兹十分重要. 例如，稍后在讨论磁本征膜\index{magnetic-eigenbrane@磁本征膜}时，我们将用到其中的一些认识.

由于我们只关注一个复结构，即 $I$，所以不会看到 $\Sigma$ 的完整超凯勒结构. 不过，这一结构有一个重要部分在复结构 $I$ 中十分清楚. 一般而言，一个超凯勒流形 $X$ 具有一个三维的实辛结构空间. 在 $X$ 的任一复结构中，三个实辛结构中的两个可以合并，成为一个全纯辛结构\index{holomorphic-symplectic-structure@全纯辛结构}的实部和虚部. 对于 $\MH$ 上的复结构 $I$，这个全纯辛结构为
\begin{equation}\label{zolk}\Omega_I=\frac{1}{\pi}\int |d ^2z|\,\Tr\,\delta\phi_z\wedge\delta A_{\bar z}. \end{equation}
这里 $-\Tr$ 是 $G$ 的李代数\index{lie-algebra@李代数} $\frak g$ 上的一个不变二次型.（我们将采用短根的长度平方为 2 的标准归一化.）在复结构 $I$ 中，超凯勒商 $\CW///\G=\MH$ 可以理解为一个复辛商，也就是说，将 $\CW$ 视为带有辛形式 $\Omega_I$ 的复辛流形，再对 $\G_\C$ 在其上的作用取辛商.（这样的辛商定义为，将复矩映射置为零，再模去复对称群；在当前情形中，复矩映射是 $\nu_I=D_{\bar z}\phi_z$，复对称群是 $\G_\C$. 在本例中，这个操作产生 $\MH$.）

我们可以写成 $\Omega_I=\omega_J+i\omega_K$，其中 $\omega_J$ 和 $\omega_K$ 是实辛形式，分别为复结构 $J$ 和 $K$ 的凯勒形式. 在几何朗兰兹中，人们主要关注辛形式为 $\omega_K$ 的 $A$-模型；我们将其简称为 $K$ 型 $A$-模型.\transnote{几何朗兰兹所用的 $A$-模型取辛形式 $\omega_K$ 而非 $\omega_J$，这是与电磁对偶相配的选择：在 $S$-对偶下 $A$-模型与 $B$-模型互换、$G$ 与 $^L\neg G$ 互换，对应关系即由此得到. 原文 §2.2.（来源：A 级，arXiv:1506.04293 §2.2；B 级 nLab “geometric Langlands correspondence” \url{https://ncatlab.org/nlab/show/geometric+Langlands+correspondence}，2026-11-19 访问；A 级，arXiv:hep-th/0604151.）}

\subsection{$\MH$ 与余切丛}
\label{hitchfib}

为了开始更详细地研究复结构 $I$ 中 $\MH$ 的几何，首先回顾，若希钦对 $(E,\varphi)$ 的底层丛 $E$ 稳定或半稳定，那么它也是稳定或半稳定的，其中 $E\to C$ 是一个全纯 $G$-丛，$\varphi\in H^0(C,K\otimes\ade)$.（反之则不成立.）因此，特别地，如果 $E$ 是一个稳定丛，那么希钦对 $(E,0)$ 总是稳定的. 这给出一个自然嵌入，将 $C$ 上稳定 $G$-丛的模空间 $\M$ 嵌入希钦模空间\index{hitchin-moduli-space@希钦模空间} $\MH$. 在复结构 $I$ 中，$\M$ 是 $\MH$ 的全纯子流形，因为它由方程 $\varphi=0$ 定义，而这个方程在复结构 $I$ 中是全纯的.（在复结构 $J$ 和 $K$ 中，$\M$ 不是全纯的. 但它对于 $\Omega_I$ 是拉格朗日的，因而对于 $\omega_J$ 和 $\omega_K$ 也是拉格朗日的.）

如果 $E$ 稳定，那么对于任意 $\varphi\in H^0(C,K\otimes\ad(E))$，希钦对 $(E,\varphi)$ 都稳定. $\M$ 在由 $E$ 代表的点处的切空间为 $H^1(C,\ad(E))$；由塞尔对偶性\index{serre-duality@塞尔对偶性}，其对偶，也就是 $\M$ 在点 $E$ 处的余切空间\index{cotangent-space@余切空间}，为 $H^0(C,K\otimes\ad(E))$. 由于 $\varphi$ 取值于这个空间，因此所有 $E$ 稳定的对 $(E,\varphi)$ 所构成的空间就是余切丛 $T^*\M$. 所以，我们实际上得到的是 $T^*\M$ 到 $\MH$ 的一个嵌入. $\MH$ 在复结构 $I$ 中的全纯辛形式 $\Omega_I$，限制在 $T^*\M$ 上，就是后者作为全纯余切丛的自然辛结构.

$T^*\M$ 在 $\MH$ 中的像并不是整个 $\MH$，因为即使 $E$ 不稳定，希钦对 $(E,\varphi)$ 仍可能稳定.\footnote{对于 $G=SU(2)$，如果 $E$ 有一个第一陈类为正的全纯线子丛 ${\cal L}$（从而 $E$ 不稳定），但 ${\cal L}$ 不是 $\varphi$-不变的（从而 $(E,\varphi)$ 稳定），就会出现这种情况. 我们将在第 \ref{distinguished} 节给出例子.} 然而，那些 $E$ 不稳定的稳定希钦对 $(E,\varphi)$，构成一个余维足够高的集合，因此在许多应用中并不重要. 去掉这个集合之后，$\MH$ 就与 $T^*\M$ 同构，并且通过忘掉 $\varphi$，自然地映到 $\M$.（文献 \cite{KW} 使用这一映射来理解 $\MH(G,C)$ 上的 $K$ 型 $A$-膜与 $\M(G,C)$ 上的 $\mathcal D$-模之间的关系.）\transnote{$T^*\mathcal M$ 的像在 $\MH$ 中是稠密开子集：$E$ 不稳定的稳定希钦对虽存在，但构成高余维集合，去掉后许多不变量不变，这正是可以用 $T^*\mathcal M$ 近似 $\MH$ 的原因；文献 [1] 由此把 $\MH$ 上的 $K$ 型 $A$-膜与 $\mathcal M$ 上的 $\mathcal D$-模联系起来. 原文 §2.3.（来源：A 级，arXiv:1506.04293 §2.3；A 级，arXiv:hep-th/0604151，2026-11-19 核对.）}

我们也可以不构造只在一般位置有定义的从 $\MH$ 到 $\M$ 的投影，而是按照稳定希钦对 $(E,\varphi)$ 的底层丛 $E$ 的全纯类型，考察 $\MH$ 的叶状结构. 这一叶状结构在 $\MH$ 的整个光滑部分上都有定义. 它的叶是中维的，并且相对于复辛结构 $\Omega_I$ 是拉格朗日的.

\subsection{希钦纤维化}
\label{hitchfibr}

\def\calP{{\mathcal P}}
通常所说的希钦纤维化\index{hitchin-fibration@希钦纤维化}，并不是把 $\varphi$ 映到 0 的映射，而是另一个映射，它把希钦对 $(E,\varphi)$ 映到 $\varphi$ 的特征多项式\index{characteristic-polynomial@特征多项式}. 对于 $G=SU(2)$，这个特征多项式就是二次微分 $w=\Tr\,\varphi^2$. 由于 $\varphi$ 全纯，$w$ 也是全纯的，因此它取值于 ${\cal V}=H^0(C,K^2)\cong \Bbb{C}^{3g-3}$. 希钦纤维化就是将 $(E,\varphi)$ 映到 $w=\Tr\,\varphi^2$ 的映射 $\ppi:\MH\to {\cal V}$.

对于任意 $G$，希钦纤维化的定义类似，只是要包括 $G$ 的全部独立卡西米尔不变量，而不只是二次卡西米尔不变量. 例如，对于 $G=SU(N)$，我们定义 $w_n=\Tr\,\varphi^n\in H^0(C,K^n)$，$n=2,\dots,N$，并令 ${\cal V}=\oplus_{n=2}^NH^0(C,K^n)$. 此时，希钦纤维化定义为将 $(E,\varphi)$ 映到 $(w_2,w_3,\dots,w_n)\in {\cal V}$. 对于任意秩为 $r$ 的李群 $G$，其李代数 $\frak g$ 上的不变多项式环是一个具有 $r$ 个生成元 ${\cal P}_i$ 的多项式代数. 这些多项式的次数 $d_i$ 满足
\begin{equation}
\label{yty}
\sum_i(2d_i-1)={\rm dim}(G).
\end{equation}
例如，对于 $G=SU(N)$，我们可以取次数为 $n$ 的多项式 $\Tr\,\varphi^n$，$n=2,\dots, N$，作为 ${\cal P}_i$，于是恒等式 (\ref{yty}) 成为 $\sum_{n=2}^N(2n-1)=N^2-1={\rm dim}(G)$. 对于任意 $G$，希钦纤维化定义为将 $(E,\varphi)$ 映到不变多项式 $\calP_i(\varphi)\in H^0(C,K^{d_i})$ 的集合. 因此，希钦纤维化的底空间为 ${\cal V}=\oplus_i H^0(C,K^{d_i})$.

由于 $\dim \,H^0(C,K^d)=(2d-1)(g-1)$，由式 (\ref{yty}) 可知，${\cal V}$ 的复维数为 $(g-1){\rm dim}(G)$，它等于 $\M$ 的维数，也等于 $\MH$ 维数的一半. 希钦纤维化 $\ppi:\MH\to {\cal V}$ 是满射，我们马上就会讨论这一点. 因此，希钦纤维化的一条一般纤维 $\FF$ 的维数也等于 $\MH$ 维数的一半：
\begin{equation}
\label{yrof}
\dim\,\FF=\dim\,{\cal V}={1\over
2}\dim\,\MH=(g-1)\dim\,G.
\end{equation}

在第 \ref{distinguished} 节中，我们将对每一点 $w\in {\cal V}$，显式构造一个稳定希钦对 $(E,\varphi)$，使它在希钦纤维化下投影到 $w$. 在此之前，我们先给出一个较为定性的论证，说明希钦纤维化是满射. 例如，取 $G=SU(2)$. 选取一个稳定的 $SU(2)$ 丛 $E$，并寻找一个希钦对 $(E,\varphi)$，使其映到 $\V$ 中由二次微分 $w$ 定义的给定点. 为此，需要找到 $\varphi\in H^0(C,K\otimes \ad(E))$，使得 $\Tr\,\varphi^2=w$. 这是关于 $3g-3$ 个未知量的 $3g-3$ 个方程，因此一般的解数为 $2^{3g-3}$. 对于其他 $G$，也可以作类似的计数.\transnote{这里的 $2^{3g-3}$ 是 Bézout 式的启发计数：$3g-3$ 个二次方程一般有 $2^{3g-3}$ 个解，它只说明 $\ppi$ 在一般点处有限、从而满射，并非纤维的真实点数；由 §2.6 的谱曲线与阿贝尔化描述，纤维实为半维数的复环面. 原文 §2.4.（来源：A 级，arXiv:1506.04293 §2.4；B 级 nLab “Hitchin fibration” \url{https://ncatlab.org/nlab/show/Hitchin+fibration}，2026-11-19 访问.）}

\subsection{完全可积性}
\label{complete}

Hitchin 的主要结果之一 \cite{HitchinDuke} 是，$\MH$ 在复结构 $I$ 中是一个完全可积哈密顿系统. 事实上，我们可以在 $\MH$ 上找到 ${1\over 2}{\rm dim}\MH$ 个函数 $H_a$，它们在复结构 $I$ 中全纯、代数独立，并且在由全纯辛形式 $\Omega_I$ 得到的泊松括号\index{poisson-bracket@泊松括号}下彼此对易.\footnote{通过选取 $\MH$ 的一个实切片，可以从这一构造中提取出一个具有实相空间和实对易哈密顿量\index{hamiltonian@哈密顿量}的通常意义下的可积系统. 一些例子和参考文献见 \cite{DonagiMarkman}.}

事实上，我们可以取 ${\cal V}$ 上的线性函数作为 $H_a$，因为 ${\cal V}$ 的维数与所需的函数数目相同. 对于 $G=SU(2)$，我们只需先在与 $H^0(C, K^2)\cong {\cal V}$ 对偶的 $(3g-3)$ 维空间 $H^1(C,T)$ 中选取一组基 $\alpha_a$，$a=1,\dots,3g-3$.（这里 $T$ 是 $C$ 的全纯切丛.）我们用取值于 $T$ 的 $(0,1)$-形式来代表 $\alpha_a$，仍将这些形式记为 $\alpha_a$，并定义
\begin{equation} \label{miro} H_a=\int_C\alpha_a \wedge
\Tr\,\varphi^2.
\end{equation}
我们断言，相对于全纯辛形式 $\Omega_I$，这些函数的泊松括号彼此为零.

一个自然的证明利用如下事实：在取超凯勒商之前，$H_a$ 的定义就在无限维空间 $\CW$ 上有意义. 利用 $\CW$ 上的辛结构 $\Omega_I$ 定义泊松括号，$H_a$ 显然是泊松对易的. 因为，根据 $\Omega_I$ 的形式 (\ref{zolk})，在这些泊松括号下，$\varphi_z$ 与自身的泊松括号为零（它与 $A_{\bar z}$ 的泊松括号当然非零）. 而 $H_a$ 仅仅是 $\varphi_z$ 的函数，不依赖于 $A_{\bar z}$.

可以将 $H_a$ 限制到满足 $0=D_{\bar z}\phi_z$ 的轨迹上；由于它们在取值于 $G_{\Bbb{C}}$ 的规范变换下不变，所以可以下降为 $\MH$ 上的全纯函数. 辛约化\index{symplectic-reduction@辛约化}的一般性质是，它将泊松对易的函数映为泊松对易的函数；这里的辛约化指的是将一个矩映射（本例中为 $\nu_I=D_{\bar z}\phi_z$）置为零，再模去相应的群（本例中为取值于 $G_{\Bbb{C}}$ 的规范变换群）. 因此，作为 $\MH$ 上的函数，$H_a$ 是泊松对易的. 它们的数目足以确立 $\MH$ 的完全可积性\index{complete-integrability@完全可积性}.

从 $G=SU(2)$ 到任意 $G$ 的推广，只需使用全部独立的规范不变多项式 ${\cal P}_i$，而不仅是 $\Tr\,\varphi^2$，来定义哈密顿量. 此时，对易哈密顿量为 $H_{a,i}=\int_C\alpha_{a,i}{\cal P}_i(\varphi)$，其中现在 $\alpha_{a,i}\in H^1(C,K^{1-d_i})$. 这些对易哈密顿量就是希钦纤维化的底空间 ${\cal V}$ 上的全部线性函数，其数目等于 $\MH$ 维数的一半.\transnote{$\alpha_{a,i}$ 的取值可由塞尔对偶性算出：$\mathcal V$ 的切空间为 $\oplus_i H^0(C,K^{d_i})$，其对偶为 $\oplus_i H^1(C,K^{1-d_i})$，因为 $K\otimes K^{-d_i}=K^{1-d_i}$；故 $\alpha_{a,i}\in H^1(C,K^{1-d_i})$ 正是与 $H^0(C,K^{d_i})$ 配对的上同调类，$H_{a,i}$ 是底空间上的线性函数. 原文 §2.5 式 (2.7).（来源：A 级，arXiv:1506.04293 §2.5；塞尔对偶性的直接计算，2026-11-19 核对.）}

在这一构造中，我们首先具体选取了用于构造对易哈密顿量的 $(0,1)$-形式 $\alpha_{a,i}$. 但是，在限制并下降到 $\MH$ 之后，所得的 $\MH$ 上的函数只依赖于 $\alpha_a$ 的上同调类. 事实上，一旦有 $\bar D\varphi=0$，从而有 $\bar\partial {\cal P}_i(\varphi)=0$，简单的分部积分便表明，$H_{a,i}$ 在 $\alpha_{a,i}\to\alpha_{a,i}+\bar\partial\epsilon_{a,i}$ 下不变.

泊松对易函数 $H_a$ 在 $\MH$ 上生成彼此对易的流，这些流在复结构 $I$ 中全纯. 此外，这些流与希钦纤维化对易，因为这些对易哈密顿量恰好是该纤维化底空间上的函数. 因而，这些流作用在希钦纤维化的各条纤维上；特别地，这些纤维容许一组极大的对易流.

复环面容许这样一组极大的对易流，因此可以猜想，至少在一般情形下，由 $H_a$ 生成的轨道是复环面. 利用希钦纤维化的“固有性”（纤维的紧致性），可以从一般论证中推出这一点；但我们将遵循文献 \cite{HitchinDuke}，利用谱曲线\index{spectral-curve@谱曲线}理论更直接地证明它.

完全可积性的一个简单而重要的推论是，希钦纤维化的纤维在全纯辛结构 $\Omega_I$ 下是拉格朗日子流形\index{lagrangian-submanifold@拉格朗日子流形}. 的确，这一纤维化的一条纤维由方程 $H_{a,i}-h_{a,i}=0$ 定义，其中 $H_{a,i}$ 是对易哈密顿量，$h_{a,i}$ 是复常数. 一般而言，由一组泊松对易函数（如本例中的 $H_{a,i}-h_{a,i}$）的消失所定义的中维子流形，是拉格朗日的.

\subsection{谱曲线}
\label{spectral}

\subsubsection{基本概念}

为了说明谱曲线的思想，\footnote{在理解以下内容时，R. Donagi 于 2004--2006 年间所作的解释给予了我极大帮助.}我们考虑 $G=SU(N)$ 的情形. 将 $E$ 视为一个秩为 $N$ 的复向量丛. 由于 $\varphi$ 取值于伴随表示，我们可以在局部把它看作由全纯 1-形式组成的一个 $N\times N$ 矩阵，并令它作用于 $E$ 的纤维. 固定一点 $\pp\in C$，并用 $\uppsi$ 表示 $E$ 在 $\pp$ 处纤维中的一个元素，便可考虑“本征值问题”
\begin{equation}
\label{udu}
\varphi(\pp)\uppsi=y\uppsi.
\end{equation}
由于 $\varphi(\pp)$ 的矩阵元取值于 $K|_\pp$，即典范丛 $K$ 在 $\pp$ 处的纤维，因此不能将 $y$ 解释为一个数. 但如果将 $y$ 解释为 $K|_\pp$ 的一个元素，这个本征值问题就有意义.

我们知道如何求一个 $N\times N$ 矩阵的本征值. 它们是特征多项式的零点. 因此，它们满足方程
\begin{equation}
\label{bludu}
\det(y-\varphi)=0.
\end{equation}
对于一般的 $\varphi$ 和 $\pp$，方程 (\ref{bludu}) 有 $N$ 个互异的根. 当 $N=2$ 时，方程可以简化. 由于 $2\times 2$ 矩阵 $\varphi$ 处处无迹，方程化为
\begin{equation}
\label{tudlu}
y^2-{1\over 2}\Tr\,\varphi^2=0.
\end{equation}

到目前为止，我们只在单个点 $\pp\in C$ 处作了说明，但显然可以对所有 $\pp$ 考虑这个方程. 当 $\pp$ 变化时，式 (\ref{bludu}) 的 $N$ 个根扫出一条代数曲线 $D$，它是 $C$ 的一个 $N$ 重覆盖. 通过“忘掉” $y$，$D$ 映到 $C$. 令 $W$ 为 $C$ 上线丛\index{line-bundle@线丛} $K$ 的全空间所构成的代数曲面. 曲线 $D$ 自然位于曲面 $W$ 中. 我们可以将 $y$ 视为纤维化 $W\to C$ 的纤维的参数. 方程 $\det(y-\varphi)=0$ 在每条纤维中选出 $N$ 个点，从而构成谱曲线 $D$，以及它的 $N$ 重覆盖映射 $\psi:D\to C$.

现在，考虑如何描述希钦纤维化 $\ppi:\MH\to {\cal V}$ 的纤维 $\FF$. 选定一条纤维 $\FF$，就意味着具体选定 $\Tr\,\varphi^2$（对于 $SU(2)$），或者选定 $\varphi$ 的特征多项式 $\det(y-\varphi)$（对于任意酉群）. 因此，一条纤维的选择确定了一条特定的谱曲线 $D$；每条纤维都与它自己的 $D$ 相关联. 如果选取的纤维足够一般，那么曲线 $D$ 是光滑且不可约的. 例如，对于 $SU(2)$，当且仅当二次微分 $\Tr\,\varphi^2$ 的零点互异时，才是如此.

这些零点就是覆盖 $D\to C$ 的分支点\index{branch-point@分支点}. $C$ 上的一个二次微分有 $4g_C-4$ 个零点. 一条亏格为 $g_C$ 的曲线，在 $4g_C-4$ 个点处分支的二重覆盖，其亏格为 $4g_C-3$. 为了验证这一点，回顾亏格为 $g_C$ 的曲线的欧拉示性数为 $2-2g_C$. 如果 $C$ 是亏格为 $g_C$ 的曲线，而 $D$ 是 $C$ 的二重覆盖，具有 $4g_C-4$ 个分支点，那么 $D$ 的欧拉示性数为 $2(2-2g_C)-(4g_C-4)$. 这等于 $g_D=4g_C-3$ 时的 $2-2g_D$，因此 $D$ 的亏格就是该值.

对于 $SU(N)$，令 $\PP(y)=\det(y-\varphi)$. 一个根为 $\lambda_i$ 的多项式 $\PP$ 的“判别式\index{discriminant@判别式}”为 $\prod_{i<j}(\lambda_i-\lambda_j)^2$. 在当前背景下，判别式是一个次数为 $N(N-1)$ 的齐次规范不变多项式 $\Delta(\varphi)$，因此是 $H^0(C,K^{N(N-1)})$ 的一个元素. 如果 $\Delta(\varphi)$ 只有单零点，那么 $D$ 光滑且不可约（更一般地，当且仅当 $\varphi$ 对每个本征值始终只有一个 Jordan 块时，才是如此）. 判别式的一个零点，就是 $C$ 上的一个点，其上方有一个 $D$ 上的点满足 $\PP(y)=\PP'(y)=0$. 方程 $\PP(y)=0$ 定义 $D$，而 $D$ 上满足 $\PP'(y)=0$ 的点恰好是覆盖 $D\to C$ 的分歧点\index{ramification-point@分歧点}. 因此，$\PP'(y)$ 的零点集在 $D$ 上定义了一个除子\index{divisor@除子} ${\cal R}$，我们称之为分歧除子\index{ramification-divisor@分歧除子}. 一般而言，它由 $2N(N-1)(g_C-1)$ 个互异的点组成（$\Delta(\varphi)\in H^0(C,K^{N(N-1)})$ 的每个零点对应一个），这意味着 $D$ 的亏格为 $g_D=g_C+(N^2-1)(g_C-1)=g_C+\dim({\cal V})=g_C+\dim(\FF)$. 如果分歧点互异，这就足以保证 $D$ 光滑.\transnote{可用黎曼–胡尔维茨公式复核：$2g_D-2=N(2g_C-2)+\deg\mathcal R$，而 $\deg\mathcal R=2N(N-1)(g_C-1)$，故 $g_D=N^2(g_C-1)+1=g_C+(N^2-1)(g_C-1)$；判别式的次数 $N(N-1)$ 正是 $\mathcal R$ 的像在 $C$ 上的次数. 原文 §2.6.1.（来源：A 级，arXiv:1506.04293 §2.6.1；A 级，Hitchin, Duke Math. J. 54 (1987) 91–114（本文文献 [21]），2026-11-19 核对.）}

现在说明一个关键事实：$\varphi$ 的本征向量在 $D$ 上构成一个线丛. 假设给定一个希钦对 $(E,\varphi)$. 对某个 $\pp\in C$，求得特征方程 (\ref{bludu}) 的一个根 $y$. 这对应于一点 $\qq\in D$，它是位于 $\pp$ 上方的 $N$ 个点 $\qq_1,\dots, \qq_N\in D$ 之一. 对每个这样的 $y$，本征向量问题 (\ref{udu}) 都有某个非零的解 $\uppsi$. 更准确地说，假设式 (\ref{bludu}) 的根互异，对于给定的 $y$，从而对于给定的 $\qq\in D$，本征向量问题具有一个一维解空间，我们将其记为 ${\cal N}_\qq$. 随着 $\qq$ 变化，${\cal N}_\qq$ 作为 $D$ 上一个复线丛 ${\cal N}$ 的纤维而变化.

当根不再互异时，线丛 ${\cal N}$ 的定义是否会失效？当曲线 $D$ 光滑时，不会. 我们只对 $G=SU(2)$ 说明这一点，一般情形与此类似. 根恰在 $\Tr\,\varphi^2=0$ 处不互异. 在 $C$ 上选取一个局部复坐标 $z$，并假定在 $z=0$ 处有 $\Tr\,\varphi^2=0$. 由于要求 $\Tr\,\varphi^2$ 具有单零点，在共轭意义下，$\varphi$ 在 $z=0$ 附近的行为为
\begin{equation}
\label{ygor}
\varphi\sim
\begin{pmatrix}0 & z\\ 1 & 0\end{pmatrix}.
\end{equation}
谱曲线在 $z=0$ 附近由下式给出：
\begin{equation}
\label{jygor}
0=\det(y-\varphi)=y^2-z.
\end{equation}
对于给定的 $z$，这一方程的解为 $y=\pm\sqrt{z}$. 不过，在 $z=0$ 附近描述谱曲线的更好方式，是将 $y$ 视为局部参数，并令 $z=y^2$. 将 $z=y^2$ 代入 $\varphi$ 的局部公式 (\ref{ygor})，本征值方程 (\ref{udu}) 成为
\begin{equation}
\label{zygor}\begin{pmatrix}0 &
y^2\\ 1 & 0\end{pmatrix}\begin{pmatrix}\psi_1
\\ \psi_2\end{pmatrix}=y\begin{pmatrix}\psi_1 \\ \psi_2\end{pmatrix}.
\end{equation}
对每一个 $y$，它都有一个一维解空间，由下式生成：
\begin{equation}
\label{calleds}
s=\begin{pmatrix}\psi_1
\\ \psi_2\end{pmatrix}=\begin{pmatrix}y \\ 1\end{pmatrix},
\end{equation}
这说明，线丛 ${\cal N}$ 的定义在谱曲线 $D$ 的简单分支点处仍然完全有效.

实际上，我们将 ${\cal N}$ 定义为 $\varphi-y$ 的核. 也可以将 $\varphi-y$ 的{\it 余核}定义为一个相关的线丛 ${\cal L}$. 只要 $\varphi$ 的本征值互异，将线性变换 $\varphi-y$ 视为从 $E\otimes K^{-1}$ 的截面到 $E$ 的截面的映射时，它就有一个一维核，因而具有一个 $(N-1)$ 维的像 $E'=\Im(\varphi-y)$. 因此，商 ${\cal L}=E/E'$ 在远离分歧点处是一维的，这就在这些点之外给出了另一个全纯线丛.

可以用上面的局部模型验证，${\cal L}$ 自然地延拓为经过分歧点的线丛. 在局部模型中，
\begin{equation}
\label{tygor}
\varphi-y=\begin{pmatrix}-y & z \\
1 & -y\end{pmatrix},
\end{equation}
因此，$\varphi-y$ 的像 $E'$ 由下式生成：
\begin{equation}
\label{defu}
u=\begin{pmatrix}-y \\ 1\end{pmatrix}.
\end{equation}
所以，商 ${\cal L}=E/E'$ 由下式生成：
\begin{equation}
\label{timely}
t=\begin{pmatrix}1\\ 0\end{pmatrix},
\end{equation}
也就是说，$u$ 和 $t$ 对每个 $y$ 都给出一组基. 这展示了 ${\cal L}$ 通过点 $z=0$ 的一个自然延拓.

${\cal L}$ 与 ${\cal N}$ 之间还存在一个重要关系. 我们将 ${\cal N}$ 定义为 $E$（更准确地说，是 $E$ 到 $D$ 的拉回\index{pullback@拉回} $\psi^*(E)$）中被 $\varphi-y$ 零化的子丛，因此有一个自然嵌入 $i:{\cal N}\to \psi^*(E)$. 同时，${\cal L}$ 是 $\psi^*(E)$ 模去 $E'=\Im(\varphi-y)$ 所得的商，因此有一个自然映射 $r:\psi^*(E)\to {\cal L}$. 复合
\begin{equation}
\label{huty}
{\cal N}\underarrow{i} \psi^*E\underarrow{r} {\cal L}
\end{equation}
给出一个全纯映射 $\theta=ri:{\cal N}\to {\cal L}$.

在远离分歧点处，可以将 $\varphi-y$ 分块对角化，使它具有一个一维核和一个余维为 1 的可逆块：
\begin{equation}\label{zelbo}\varphi-y=\begin{pmatrix} 0 & 0\cr 0&*\end{pmatrix}.\end{equation}
因此，核与余核都由上方的元素（零模）生成，而映射 $\theta:\cal N\to \cal L$ 是一个同构.

为了考察分歧点处的情形，再次使用 $N=2$ 的局部模型. ${\cal N}$ 由式 (\ref{calleds}) 中记作 $s$ 的向量生成. 为了计算 $\theta(s)$，只需模去式 (\ref{defu}) 所定义的 $u$，再将结果表示成式 (\ref{timely}) 中 ${\cal L}$ 的生成元 $t$ 的倍数. 有 $s=u+2yt$，即
\begin{equation}
\label{burble}
s = 2yt ~{\rm mod}~ u.
\end{equation}
因此，${\cal N}$ 的一个生成元 $s$ 映到 ${\cal L}$ 的一个生成元的 $2y$ 倍. 由于 $y$ 在分歧除子 ${\cal R}$ 处有一个单零点，结果是
\begin{equation}
\label{troygo}
{\cal L}\cong {\cal N}({\cal R}),
\end{equation}
换言之，${\cal L}$ 的一个截面等价于 ${\cal N}$ 的一个截面，后者允许在分歧除子处有极点.\transnote{${\cal N}(\mathcal R)$ 是允许沿分歧除子 $\mathcal R$ 取极点的扭化丛，即 ${\cal N}\otimes\mathcal O_D(\mathcal R)$；它与式 (2.22) 的写法一致：$y$ 是 $\psi^*(K)$ 的截面、$\PP'(y)$ 关于 $y$ 为 $N-1$ 次，零点除子恰为 $\mathcal R$，于是 $\mathcal R$ 对应 $\psi^*(K^{N-1})$. 原文 §2.6.1 式 (2.21)–(2.22).（来源：A 级，arXiv:1506.04293 §2.6.1；A 级，Hitchin, Duke Math. J. 54 (1987)（本文文献 [21]），2026-11-19 核对.）}

如上所述，分歧除子恰好是多项式 $\PP'(y)$ 在 $D$ 上的零点集，而它是 $\psi^*(K^{N-1})$ 的一个截面（因为 $y$ 是 $\psi^*(K)$ 的截面，且 $\PP'(y)$ 关于 $y$ 的次数为 $N-1$）. 因而也可以写成
\begin{equation}
\label{roygo} {\cal L}\cong {\cal N}\otimes \psi^*(K^{N-1}).
\end{equation}

\subsubsection{阿贝尔化}

现在，关键的认识是，可以从线丛 ${\cal L}\to D$ 重建丛 $E\to C$（以及希格斯场 $\varphi$）. 这是一项进展，因为 $E$ 是一个 $G$-丛，其中 $G$ 是{\it 非阿贝尔}的，而线丛则对应于将 $G$ 换成阿贝尔群 $U(1)$ 时的 $G$-丛. 为这种阿贝尔化\index{abelianization@阿贝尔化}所付出的代价是，我们必须在谱覆盖\index{spectral-cover@谱覆盖} $D$ 上工作，而不再是在曲线 $C$ 上工作.\transnote{这正是本节标题所说的“阿贝尔化”：把 $C$ 上非阿贝尔的 $G$-丛 $E$ 换成谱覆盖 $D$ 上的线丛 ${\cal L}$，代价是工作曲线从 $C$ 换成 $D$——$D\to C$ 是 $N$ 重覆盖，亏格为 $N^2(g_C-1)+1$，远大于 $g_C$；希格斯场则由 $\varphi=\psi_*(y)$ 一并恢复. 原文 §2.6.2 式 (2.25)、(2.29).（来源：A 级，arXiv:1506.04293 §2.6.2；A 级，arXiv:hep-th/0604151，2026-11-19 核对.）}

首先，假设点 $\pp\in C$ 不是纤维化 $\psi:D\to C$ 的分支点. 那么，$E$ 在 $\pp$ 处的纤维，记为 $E_\pp$，可以分解为 $\varphi$ 的 $N$ 个一维本征空间之和. 这些本征空间就是上面称为 ${\cal N}$ 的线丛，在位于 $\pp$ 上方的点 $\qq_i$ 处的纤维. 因此，只要 $\pp$ 不是分支点，就有
\begin{equation}
\label{ytot}
E_\pp=\oplus_{i=1}^N{\cal N}_{\qq_i}.
\end{equation}
这无疑表明，在远离分支点处，可以由 ${\cal N}$ 恢复 $E$. 由于 ${\cal N}$ 与 ${\cal L}$ 在远离分支点处自然同构，所以在分支点之外，同样可以写成
\begin{equation}
\label{bytot}
E_\pp=\oplus_{i=1}^N{\cal L}_{\qq_i}
\end{equation}
事实证明，这种描述延拓到分支点处更加简单.

将这个公式延拓到分支点处，需要用到代数几何中的另一个概念，即“推前\index{pushforward@推前}”. 对任意映射 $\psi:D\to C$ 和 $D$ 上的任意层，都可以定义推前. 我们考虑的层是线丛或向量丛的截面层\index{sheaf@截面层}，在记号上不区分一个丛和相应的截面层. 如果 ${\cal L}$ 是 $D$ 上的一个线丛，其推前 $\psi_*({\cal L})$ 就是 $C$ 上按如下方式定义的层：$\psi_*({\cal L})$ 在一个充分小的开集 $U\subset C$ 上的截面，与 ${\cal L}$ 在 $\psi^{-1}(U)$ 上的截面相同.

我们断言，可以由 ${\cal L}$ 按如下方式重建 $E$：
\begin{equation}
\label{hytot}E=\psi_*({\cal L}).
\end{equation}
在远离分歧点处，这只不过是式 (\ref{bytot}) 的一种较为精巧的重述. 我们再次利用局部模型，看看在分歧点处会发生什么. 虽然结果是一般的，这里只讨论 $N=2$.

${\cal L}$ 的截面与 $\psi^*E$ 的截面相同，只是必须模去 ${\rm Im}(\varphi-y)$ 的生成元 $u$. 非正式地说，$\psi_*(\psi^*E)$ 的截面就是 $E$ 的截面，但允许它依赖于 $y$. 因此，$\psi_*({\cal L})$ 的截面是满足如下条件的 $E$ 的截面：{\it (i)} 它可以依赖于 $y$；{\it (ii)} 如果它是 $u$ 的倍数，就视为平凡的.

$E$ 的一个截面具有如下形式：
\begin{equation}
\label{tyro}
\begin{pmatrix}a(z)\\
b(z)\end{pmatrix},
\end{equation}
其中 $a(z)$、$b(z)$ 是全纯函数. $\psi_*({\cal L})$ 的一个截面可以写成
\begin{equation}
\label{wytot}
\begin{pmatrix}A(z)+yC(z)\\ B(z)+yD(z)\end{pmatrix}~{\rm
mod}~ \left(G(z)+yH(z)\right)\begin{pmatrix}-y \\ 1\end{pmatrix},
\end{equation}
这里允许依赖于 $y$，并模去 $u$. 对于 $G=-C$、$H=D$，有
\begin{equation}
\label{qytot}
(G(z)+yH(z))\begin{pmatrix}-y \\
1\end{pmatrix} =y\begin{pmatrix}C \\ D\end{pmatrix}
+\begin{pmatrix}-zD\\ -C\end{pmatrix},
\end{equation}
这表明，式 (\ref{wytot}) 中的等价关系足以唯一地令 $C=D=0$，从而使式 (\ref{wytot}) 等价于式 (\ref{tyro})，并说明陈述 $E=\psi_*({\cal L})$ 在分歧点处仍然成立.

类似地，可以通过下式恢复希格斯场 $\varphi$：
\begin{equation}
\label{vyro} \varphi=\psi_*(y).
\end{equation}
这个公式的意思是，如果 $f$ 是 ${\cal L}\to D$ 的一个截面，而 $s=\psi_*(f)$ 是 $E\to C$ 的相应截面，那么 $\varphi(s)=\psi_*(yf)$，这一条件足以确定 $\varphi$. 为了说明这一点，注意 ${\cal L}$ 的一个截面 $f$，就是 $\psi^*(E)$ 的一个截面 $\hat f$ 模去如下等价关系：对于任意 $\chi$，令 $(\varphi-y)\chi$ 为零. 在这个等价关系下，$y\hat f=\varphi\hat f$. 由于这对所有 $\hat f$ 都成立，将其推到 $C$ 上，就得到关系式 (\ref{vyro}).

\subsubsection{会出现哪些线丛？}
\label{whichline}

因此，每个希钦对 $(E,\varphi)$ 都来自谱覆盖 $D$ 上的一个线丛 ${\cal L}$. 但哪些线丛会以这种方式出现呢？我们先考察 ${\cal L}$ 的次数，再考察它的模参数.

首先计算推前 $\psi_*({\cal O})$，其中 ${\cal O}$ 是 $D$ 上的平凡线丛. ${\cal O}\to D$ 的一个局部全纯截面 $s$ 可以按 $y$ 的幂展开：
\begin{equation}
\label{rofo}
s=\alpha_0+y\alpha_1+y^2\alpha_2+\dots+y^{N-1}\alpha_{N-1}.
\end{equation}
由于方程 $\det(y-\varphi)=0$，展开式在 $y^{N-1}$ 处终止. 由于 $y$ 是 $\psi^*(K)$ 的截面，每个 $\alpha_i$ 都是 $C$ 上线丛 $K^{-i}$ 的截面. 因此
\begin{equation}
\label{jofo}\psi_*({\cal O})={\cal O}\oplus K^{-1}\oplus
K^{-2}\oplus\dots\oplus K^{-(N-1)}.
\end{equation}
为便于后面使用，我们也可以容易地计算当 ${\cal L}=\psi^*({\cal L}_0)$ 是 $C$ 上一个线丛 ${\cal L}_0$ 的拉回时的 $\psi_*({\cal L})$. 此时，${\cal L}$ 的局部截面 $s$ 可以像式 (\ref{rofo}) 一样展开，但现在 $\alpha_i$ 是 ${\cal L}_0\otimes K^{-i}$ 的截面. 因而
\begin{equation}
\label{xofo} \psi_*(\psi^*({\cal L}_0))={\cal
L}_0\otimes\left({\cal O}\oplus K^{-1}\oplus K^{-2}\oplus
\dots\oplus K^{-(N-1)}\right).
\end{equation}

由于 $K^{-i}$ 的次数为 $-2i(g_C-1)$，式 (\ref{jofo}) 意味着 $\psi_*({\cal O})$ 的第一陈类为 $d=-N(N-1)(g_C-1)$. 如果改为取次数为 $c$ 的 ${\cal L}\to D$，那么 $\psi_*({\cal L})$ 就是 $C$ 上第一陈类为 $c+d$ 的向量丛.\footnote{说明这一点的一种方法，是使用黎曼-罗赫公式\index{riemann-roch@黎曼-罗赫公式}，以及 ${\cal L}\to D$ 的全纯欧拉示性数与 $\psi_*({\cal L})\to C$ 的全纯欧拉示性数相等这一事实. 如果 ${\cal L}\to D$ 的第一陈类增加 1，全纯欧拉示性数就增加 1，因此 $\psi_*({\cal L})\to C$ 的第一陈类也必须增加 1.} 要得到一个 $SU(N)$ 丛，第一陈类应当为零，因此 ${\cal L}$ 的次数必须为
\begin{equation}
\label{ywo} c_0=N(N-1)(g_C-1).
\end{equation}

\def\Jac{{\rm Jac}}
因此，与一个 $SU(N)$ 希钦对相关联的线丛 ${\cal L}$ 具有这个次数，并代表 ${\rm Pic}_{c_0}(D)$ 中的一个点，即 $D$ 的皮卡群\index{picard-group@皮卡群}中参数化次数为 $c_0$ 的线丛的分支.

皮卡群每个分支的维数都等于 $D$ 的亏格 $g_D$，而 $g_D=g_C+\dim(\FF)$. 因此，线丛 ${\cal L}$ 所依赖的连续参数的数目，比希钦纤维化的纤维 $\FF$ 的维数多 $g_C$. 为了参数化这条纤维，必须对 ${\cal L}$ 施加 $g_C$ 个条件.

这些条件应当是什么，是很清楚的. 关于 ${\cal L}$ 次数的条件 (\ref{ywo}) 保证行列式线丛\index{determinant-line-bundle@行列式线丛} $\det(E)$ 的次数为零，因此在拓扑上平凡. 要得到一个 $SU(N)$ 希钦对 $(E,\varphi)$，$\det(E)$ 还必须在全纯意义下平凡. {\it 先验地}，$\det(E)$ 取值于 $C$ 的雅可比簇\index{jacobian@雅可比簇} $\Jac(C)$，其维数为 $g_C$. 因而，要求 $\det(E)$ 全纯平凡，就施加了 $g_C$ 个条件，将维数减到正确的数值. 若用 $\Lambda$ 表示将 $D$ 上的线丛 ${\cal L}$ 映到 $C$ 上的线丛 $\Lambda({\cal L})=\det\psi_*({\cal L})$ 的映射，那么 $SU(N)$ 的希钦纤维就由满足 $\Lambda({\cal L})\cong {\cal O}$ 的线丛组成. 换言之，这条纤维为
\begin{equation}
\label{yudo}
\FF_{SU(N)}=\Lambda^{-1}({\cal O}).
\end{equation}

在远离分支点处，可以更简单地表达 ${\cal L}$ 所满足的条件. 如果 $\pp$ 不是分支点，那么 $E_\pp$ 具有分解式 (\ref{bytot})，由此有 $\det(E)_\pp=\otimes_{i=1}^N{\cal L}_{\qq_i}$. 因此，在远离分支点处，存在一个全纯变化的同构
\begin{equation}
\label{milto}
\otimes_{i=1}^N{\cal L}_{\qq_i}\cong \Bbb{C}.
\end{equation}

一般而言，复环面之间一个全纯映射的核，在本例中即 $\Lambda:{\rm Pic}_{c_0}(D)\to {\rm Jac}(C)$ 的核，是一个复环面或若干环面的并. 我们将在第 \ref{topview} 节证明，$\FF_{SU(N)}$ 连通，因此它实际上是一个复环面. 它是否以一种自然的方式成为阿贝尔簇\index{abelian-variety@阿贝尔簇}呢？要把它呈现为阿贝尔簇，就必须以自然的方式在 $\FF_{SU(N)}$ 中选取一点，也就是说，必须选取一个具有适当次数且位于 $\Lambda$ 的核中的线丛 ${\cal L}\to D$. 为此，选取 $C$ 的典范丛的一个平方根 $K^{1/2}$，并令 ${\cal L}_0=\psi^*(K^{(N-1)/2})$. 它的确具有适当的次数（$K^{(N-1)/2}$ 的次数为 $(g_C-1)(N-1)$，而由于 ${\cal L}_0$ 是经由一个次数为 $N$ 的映射所作的拉回，它的次数为原来的 $N$ 倍）. ${\cal L}_0$ 也确实位于映射 $\Lambda$ 的核中. 事实上，由式 (\ref{xofo}) 可知，在这种情况下有 $\psi_*({\cal L})=K^{(N-1)/2}\otimes\left({\cal O}\oplus K^{-1}\oplus K^{-2}\oplus \dots\oplus K^{-(N-1)}\right)$，由此得到 $\det(\psi_*({\cal L}_0))\cong {\cal O}$.

如果 $N$ 为奇数，$(N-1)/2$ 就是整数，而这一构造实际上并不依赖于 $K$ 的平方根的选择. 因此，对于奇数 $N$，我们已经将 $\FF_{SU(N)}$ 呈现为一个阿贝尔簇，并为希钦纤维化的每条纤维给出了一个特选代表 $(E,\varphi)$. 对于偶数 $N$，情形并非完全如此，因为 $K^{1/2}$ 的选择确实有影响. 当 $N$ 为偶数时，我们最多只能构造这样一对 $(E,\varphi)$：除去可能与一个二阶平坦线丛\index{flat-line-bundle@平坦线丛}作张量积的自由度之外，它是自然的. 这一区别由文献 \cite{donga} 指出，并推广到其他 $G$. 也可参见下文第 \ref{distinguished} 节末尾的讨论.\transnote{取基点后复环面才成为阿贝尔簇：$\Lambda({\cal L})=\det\psi_*({\cal L})$ 把 $\mathrm{Pic}_{c_0}(D)$ 映到 $\mathrm{Jac}(C)$，其核 $\mathcal F_{SU(N)}$ 的维数为 $g_D-g_C=\dim\mathcal F$，正是一个“Prym 型”对象（推前为平凡线丛的线丛构成的空间）. ${\cal L}_0=\psi^*(K^{(N-1)/2})$ 给出核中一个自然点：奇 $N$ 时它与 $K^{1/2}$ 的选取无关，偶 $N$ 时只差一个二阶平坦线丛的自由度，这一差别由文献 [23] 指出并推广到其他 $G$. 原文 §2.6.3.（来源：A 级，arXiv:1506.04293 §2.6.3；B 级 nLab “Prym variety” \url{https://ncatlab.org/nlab/show/Prym+variety}，2026-11-19 访问；A 级，arXiv:math/0005132（本文文献 [23]）.）}
